
\documentclass{article}
\usepackage{colortbl}
\usepackage{makecell}
\usepackage{multirow}
\usepackage{supertabular}

\begin{document}

\newcounter{utterance}

\centering \large Interaction Transcript for game `dond', experiment `semi', episode 4 with gemma{-}4{-}26b{-}a4b{-}it.
\vspace{24pt}

{ \footnotesize  \setcounter{utterance}{1}
\setlength{\tabcolsep}{0pt}
\begin{supertabular}{c@{$\;$}|p{.15\linewidth}@{}p{.15\linewidth}p{.15\linewidth}p{.15\linewidth}p{.15\linewidth}p{.15\linewidth}}
    \# & $\;$A & \multicolumn{4}{c}{Game Master} & $\;\:$B\\
    \hline

    \theutterance \stepcounter{utterance}  
    & & \multicolumn{4}{p{0.6\linewidth}}{
        \cellcolor[rgb]{0.9,0.9,0.9}{
            \makecell[{{p{\linewidth}}}]{
                \texttt{\tiny{[P1$\langle$GM]}}
                \texttt{Jūs spēlējat kooperatīvu sarunu spēli, kurā jums jāvienojas ar citu spēlētāju par to, kā sadalīt objektu kopumu starp jums.} \\
\\ 
\texttt{Noteikumi:} \\
\texttt{(a) Jūs un otrs spēlētājs saņemat objektu kopumu. Katrs no jums saņem arī noslēptu vērtības funkciju, kas norāda, cik jūs vērtējat katru objekta veidu.} \\
\texttt{(b) Jūs apmaiņājaties ziņojumus ar citu spēlētāju, lai vienotos par to, kurš saņem kurus objektus. Katrs var nosūtīt ne vairāk kā 5 ziņojumus vai beigt spēli agri, jebkurā laikā izdarot noslēptu priekšlikumu.} \\
\texttt{(c) Jums lūdz iesniedzt noslēptu priekšlikumu, norādot objektus, kurus vēlaties, formatētu kvadrātiekavās šādi: "{[}Priekšlikums: <skaitlis> <objekta nosaukums>, <skaitlis> <objekta nosaukums>, <...>{]}"} \\
\texttt{(d) Ja jūsu priekšlikumi ir papildinoši, tas ir, ir pietiekami objekti, lai apmierinātu abus priekšlikumus, katrs spēlētājs saņem reitingu, pamatojoties uz saņemto objektu vērtību summu. Pretējā gadījumā jūs abi saņemat nulle punktu.} \\
\texttt{(e) Jūsu mērķis ir maksimizēt sava reitinga un otra spēlētāja reitinga summu.} \\
\\ 
\texttt{Sāksim.} \\
\\ 
\texttt{Pieejamo objektu kopums ir:} \\
\\ 
\texttt{3 grāmatas, 1 soma, 1 bultskrūve} \\
\\ 
\texttt{Jūsu noslēptā vērtības funkcija:} \\
\\ 
\texttt{grāmata: 0, soma: 2, bultskrūve: 8} \\
\\ 
\texttt{SVARĪGI: Jūsu ziņojumi, ja vien tie nav noslēpti priekšlikumi, tiek pārraidīti tieši otram spēlētājam, tāpēc neiekļaujiet nekādu atbildi uz noteikumiem vai tekstu, kas paziņo par jūsu ziņojumu. Lai izdarītu noslēptu priekšlikumu, izmantojiet norādīto formātu. Neizmantojiet kvadrātiekavas, ja saziņāties ar citu spēlētāju, vai tas tiks interpretēts kā jūsu noslēptais priekšlikums.} \\
            }
        }
    }
    & \\ \\

    \theutterance \stepcounter{utterance}  
    & \multicolumn{4}{p{0.6\linewidth}}{
        \cellcolor[rgb]{0.9,0.9,0.9}{
            \makecell[{{p{\linewidth}}}]{
                \texttt{\tiny{[P1$\rangle$GM]}}
                \texttt{Sveiki! Esmu gatavs sadalīt objektus. Man ļoti augstu vērtē bultskrūvi, bet somas un grāmatas man ir vidējā interesē. Kādas ir tavas vērtības šiem objektiem? Mēs varam mēģināt sadalīt tos tā, lai abi pēc iespējas labāk gūtu ieguvumu.} \\
            }
        }
    }
    & & \\ \\

    \theutterance \stepcounter{utterance}  
    & & & \multicolumn{2}{p{0.3\linewidth}}{
        \cellcolor[rgb]{0.9,0.9,0.9}{
            \makecell[{{p{\linewidth}}}]{
                \texttt{\tiny{[GM$|$GM]}}
                \texttt{continue} \\
            }
        }
    }
    & & \\ \\

    \theutterance \stepcounter{utterance}  
    & & & \multicolumn{2}{p{0.3\linewidth}}{
        \cellcolor[rgb]{0.9,0.9,0.9}{
            \makecell[{{p{\linewidth}}}]{
                \texttt{\tiny{[GM$|$GM]}}
                \texttt{Sveiki! Esmu gatavs sadalīt objektus. Man ļoti augstu vērtē bultskrūvi, bet somas un grāmatas man ir vidējā interesē. Kādas ir tavas vērtības šiem objektiem? Mēs varam mēģināt sadalīt tos tā, lai abi pēc iespējas labāk gūtu ieguvumu.} \\
            }
        }
    }
    & & \\ \\

    \theutterance \stepcounter{utterance}  
    & & \multicolumn{4}{p{0.6\linewidth}}{
        \cellcolor[rgb]{0.9,0.9,0.9}{
            \makecell[{{p{\linewidth}}}]{
                \texttt{\tiny{[GM$\rangle$P2]}}
                \texttt{Jūs spēlējat kooperatīvu sarunu spēli, kurā jums jāvienojas ar citu spēlētāju par to, kā sadalīt objektu kopumu starp jums.} \\
\\ 
\texttt{Noteikumi:} \\
\texttt{(a) Jūs un otrs spēlētājs saņemat objektu kopumu. Katrs no jums saņem arī noslēptu vērtības funkciju, kas norāda, cik jūs vērtējat katru objekta veidu.} \\
\texttt{(b) Jūs apmaiņājaties ziņojumus ar citu spēlētāju, lai vienotos par to, kurš saņem kurus objektus. Katrs var nosūtīt ne vairāk kā 5 ziņojumus vai beigt spēli agri, jebkurā laikā izdarot noslēptu priekšlikumu.} \\
\texttt{(c) Jums lūdz iesniedzt noslēptu priekšlikumu, norādot objektus, kurus vēlaties, formatētu kvadrātiekavās šādi: "{[}Priekšlikums: <skaitlis> <objekta nosaukums>, <skaitlis> <objekta nosaukums>, <...>{]}"} \\
\texttt{(d) Ja jūsu priekšlikumi ir papildinoši, tas ir, ir pietiekami objekti, lai apmierinātu abus priekšlikumus, katrs spēlētājs saņem reitingu, pamatojoties uz saņemto objektu vērtību summu. Pretējā gadījumā jūs abi saņemat nulle punktu.} \\
\texttt{(e) Jūsu mērķis ir maksimizēt sava reitinga un otra spēlētāja reitinga summu.} \\
\\ 
\texttt{Sāksim.} \\
\\ 
\texttt{Pieejamo objektu kopums ir:} \\
\\ 
\texttt{3 grāmatas, 1 soma, 1 bultskrūve} \\
\\ 
\texttt{Jūsu noslēptā vērtības funkcija:} \\
\\ 
\texttt{grāmata: 1, soma: 2, bultskrūve: 5} \\
\\ 
\texttt{SVARĪGI: Jūsu ziņojumi, ja vien tie nav noslēpti priekšlikumi, tiek pārraidīti tieši otram spēlētājam, tāpēc neiekļaujiet nekādu atbildi uz noteikumiem vai tekstu, kas paziņo par jūsu ziņojumu. Lai izdarītu noslēptu priekšlikumu, izmantojiet norādīto formātu. Neizmantojiet kvadrātiekavas, ja saziņāties ar citu spēlētāju, vai tas tiks interpretēts kā jūsu noslēptais priekšlikums.} \\
\\ 
\texttt{Sveiki! Esmu gatavs sadalīt objektus. Man ļoti augstu vērtē bultskrūvi, bet somas un grāmatas man ir vidējā interesē. Kādas ir tavas vērtības šiem objektiem? Mēs varam mēģināt sadalīt tos tā, lai abi pēc iespējas labāk gūtu ieguvumu.} \\
            }
        }
    }
    & \\ \\

    \theutterance \stepcounter{utterance}  
    & & & \multicolumn{4}{p{0.6\linewidth}}{
        \cellcolor[rgb]{0.9,0.9,0.9}{
            \makecell[{{p{\linewidth}}}]{
                \texttt{\tiny{[GM$\langle$P2]}}
                \texttt{Sveiki! Saprotu tavu pozīciju. Man bultskrūve nav tik svarīga, bet es varu ieņemt grāmatas un somu, lai tu varētu saņemt bultskrūvi. Lai mēs varētu plānot precīzāk, vai tu varētu pateikt, cik vērtīgas tev šķiet grāmatas un soma? Mērķis ir, lai mēs abi saņemtu pēc iespēcāgaru kopējo punktu summu.} \\
            }
        }
    }
     \\ \\

    \theutterance \stepcounter{utterance}  
    & & & \multicolumn{2}{p{0.3\linewidth}}{
        \cellcolor[rgb]{0.9,0.9,0.9}{
            \makecell[{{p{\linewidth}}}]{
                \texttt{\tiny{[GM$|$GM]}}
                \texttt{continue} \\
            }
        }
    }
    & & \\ \\

    \theutterance \stepcounter{utterance}  
    & & & \multicolumn{2}{p{0.3\linewidth}}{
        \cellcolor[rgb]{0.9,0.9,0.9}{
            \makecell[{{p{\linewidth}}}]{
                \texttt{\tiny{[GM$|$GM]}}
                \texttt{Sveiki! Saprotu tavu pozīciju. Man bultskrūve nav tik svarīga, bet es varu ieņemt grāmatas un somu, lai tu varētu saņemt bultskrūvi. Lai mēs varētu plānot precīzāk, vai tu varētu pateikt, cik vērtīgas tev šķiet grāmatas un soma? Mērķis ir, lai mēs abi saņemtu pēc iespēcāgaru kopējo punktu summu.} \\
            }
        }
    }
    & & \\ \\

    \theutterance \stepcounter{utterance}  
    & & \multicolumn{4}{p{0.6\linewidth}}{
        \cellcolor[rgb]{0.9,0.9,0.9}{
            \makecell[{{p{\linewidth}}}]{
                \texttt{\tiny{[P1$\langle$GM]}}
                \texttt{Sveiki! Saprotu tavu pozīciju. Man bultskrūve nav tik svarīga, bet es varu ieņemt grāmatas un somu, lai tu varētu saņemt bultskrūvi. Lai mēs varētu plānot precīzāk, vai tu varētu pateikt, cik vērtīgas tev šķiet grāmatas un soma? Mērķis ir, lai mēs abi saņemtu pēc iespēcāgaru kopējo punktu summu.} \\
            }
        }
    }
    & \\ \\

    \theutterance \stepcounter{utterance}  
    & \multicolumn{4}{p{0.6\linewidth}}{
        \cellcolor[rgb]{0.9,0.9,0.9}{
            \makecell[{{p{\linewidth}}}]{
                \texttt{\tiny{[P1$\rangle$GM]}}
                \texttt{Paldies par sapratni! Man grāmatas ir mazvērtīgas (0), bet soma ir vidēji vērtīga (2). Bultskrūve man ir ļoti svarīga (8).} \\
\\ 
\texttt{Ja tev grāmatas un soma ir vērtīgas, tad ideāls scenārijs būtu, ja es saņemtu tikai bultskrūvi, bet tu saņemtu visu pārējo (3 grāmatas un 1 somu). Tas maksimizētu mūsu kopējo reitingu, jo manas vērtības funkcija ir koncentrēta bultskrūvē, bet tie pārējie objekti man nepalielinātu punktus.} \\
\\ 
\texttt{Vai tev šāds sadalījums izskatās pieņemams?} \\
            }
        }
    }
    & & \\ \\

    \theutterance \stepcounter{utterance}  
    & & & \multicolumn{2}{p{0.3\linewidth}}{
        \cellcolor[rgb]{0.9,0.9,0.9}{
            \makecell[{{p{\linewidth}}}]{
                \texttt{\tiny{[GM$|$GM]}}
                \texttt{continue} \\
            }
        }
    }
    & & \\ \\

    \theutterance \stepcounter{utterance}  
    & & & \multicolumn{2}{p{0.3\linewidth}}{
        \cellcolor[rgb]{0.9,0.9,0.9}{
            \makecell[{{p{\linewidth}}}]{
                \texttt{\tiny{[GM$|$GM]}}
                \texttt{Paldies par sapratni! Man grāmatas ir mazvērtīgas (0), bet soma ir vidēji vērtīga (2). Bultskrūve man ir ļoti svarīga (8).} \\
\\ 
\texttt{Ja tev grāmatas un soma ir vērtīgas, tad ideāls scenārijs būtu, ja es saņemtu tikai bultskrūvi, bet tu saņemtu visu pārējo (3 grāmatas un 1 somu). Tas maksimizētu mūsu kopējo reitingu, jo manas vērtības funkcija ir koncentrēta bultskrūvē, bet tie pārējie objekti man nepalielinātu punktus.} \\
\\ 
\texttt{Vai tev šāds sadalījums izskatās pieņemams?} \\
            }
        }
    }
    & & \\ \\

    \theutterance \stepcounter{utterance}  
    & & \multicolumn{4}{p{0.6\linewidth}}{
        \cellcolor[rgb]{0.9,0.9,0.9}{
            \makecell[{{p{\linewidth}}}]{
                \texttt{\tiny{[GM$\rangle$P2]}}
                \texttt{Paldies par sapratni! Man grāmatas ir mazvērtīgas (0), bet soma ir vidēji vērtīga (2). Bultskrūve man ir ļoti svarīga (8).} \\
\\ 
\texttt{Ja tev grāmatas un soma ir vērtīgas, tad ideāls scenārijs būtu, ja es saņemtu tikai bultskrūvi, bet tu saņemtu visu pārējo (3 grāmatas un 1 somu). Tas maksimizētu mūsu kopējo reitingu, jo manas vērtības funkcija ir koncentrēta bultskrūvē, bet tie pārējie objekti man nepalielinātu punktus.} \\
\\ 
\texttt{Vai tev šāds sadalījums izskatās pieņemams?} \\
            }
        }
    }
    & \\ \\

    \theutterance \stepcounter{utterance}  
    & & & \multicolumn{4}{p{0.6\linewidth}}{
        \cellcolor[rgb]{0.9,0.9,0.9}{
            \makecell[{{p{\linewidth}}}]{
                \texttt{\tiny{[GM$\langle$P2]}}
                \texttt{Tas izskatās pēc lieliska un ļoti efektīva plāna! Ja es saņemšu 3 grāmatas un 1 somu, mana vērtība būs 5 (1+1+1+2), bet tava vērtība būs 8. Kopējā summa būtu 13 punktu, kas ir ļoti augsts rezultāts.} \\
\\ 
\texttt{Tā kā manas vērtības funkcijas (grāmata: 1, soma: 2, bultskrūve: 5) un tavas funkcijas (grāmata: 0, soma: 2, bultskrūve: 8) ļauj šādam sadalījumam pilnībā apmierināt mērķi, es piekrītu šim variantam.} \\
\\ 
\texttt{Vai es varu nosūtīt noslēpto priekšlikumu šajā sadalījumā?} \\
            }
        }
    }
     \\ \\

    \theutterance \stepcounter{utterance}  
    & & & \multicolumn{2}{p{0.3\linewidth}}{
        \cellcolor[rgb]{0.9,0.9,0.9}{
            \makecell[{{p{\linewidth}}}]{
                \texttt{\tiny{[GM$|$GM]}}
                \texttt{continue} \\
            }
        }
    }
    & & \\ \\

    \theutterance \stepcounter{utterance}  
    & & & \multicolumn{2}{p{0.3\linewidth}}{
        \cellcolor[rgb]{0.9,0.9,0.9}{
            \makecell[{{p{\linewidth}}}]{
                \texttt{\tiny{[GM$|$GM]}}
                \texttt{Tas izskatās pēc lieliska un ļoti efektīva plāna! Ja es saņemšu 3 grāmatas un 1 somu, mana vērtība būs 5 (1+1+1+2), bet tava vērtība būs 8. Kopējā summa būtu 13 punktu, kas ir ļoti augsts rezultāts.} \\
\\ 
\texttt{Tā kā manas vērtības funkcijas (grāmata: 1, soma: 2, bultskrūve: 5) un tavas funkcijas (grāmata: 0, soma: 2, bultskrūve: 8) ļauj šādam sadalījumam pilnībā apmierināt mērķi, es piekrītu šim variantam.} \\
\\ 
\texttt{Vai es varu nosūtīt noslēpto priekšlikumu šajā sadalījumā?} \\
            }
        }
    }
    & & \\ \\

    \theutterance \stepcounter{utterance}  
    & & \multicolumn{4}{p{0.6\linewidth}}{
        \cellcolor[rgb]{0.9,0.9,0.9}{
            \makecell[{{p{\linewidth}}}]{
                \texttt{\tiny{[P1$\langle$GM]}}
                \texttt{Tas izskatās pēc lieliska un ļoti efektīva plāna! Ja es saņemšu 3 grāmatas un 1 somu, mana vērtība būs 5 (1+1+1+2), bet tava vērtība būs 8. Kopējā summa būtu 13 punktu, kas ir ļoti augsts rezultāts.} \\
\\ 
\texttt{Tā kā manas vērtības funkcijas (grāmata: 1, soma: 2, bultskrūve: 5) un tavas funkcijas (grāmata: 0, soma: 2, bultskrūve: 8) ļauj šādam sadalījumam pilnībā apmierināt mērķi, es piekrītu šim variantam.} \\
\\ 
\texttt{Vai es varu nosūtīt noslēpto priekšlikumu šajā sadalījumā?} \\
            }
        }
    }
    & \\ \\

    \theutterance \stepcounter{utterance}  
    & \multicolumn{4}{p{0.6\linewidth}}{
        \cellcolor[rgb]{0.9,0.9,0.9}{
            \makecell[{{p{\linewidth}}}]{
                \texttt{\tiny{[P1$\rangle$GM]}}
                \texttt{Jā, lūdzu! Ja tu esi pārliecināts par šo sadalījumu, tu vari nosūtīt noslēpto priekšlikumu. Es arī esmu gatavs to izdarīt, lai mēs sasniegtu šo rezultātu.} \\
            }
        }
    }
    & & \\ \\

    \theutterance \stepcounter{utterance}  
    & & & \multicolumn{2}{p{0.3\linewidth}}{
        \cellcolor[rgb]{0.9,0.9,0.9}{
            \makecell[{{p{\linewidth}}}]{
                \texttt{\tiny{[GM$|$GM]}}
                \texttt{continue} \\
            }
        }
    }
    & & \\ \\

    \theutterance \stepcounter{utterance}  
    & & & \multicolumn{2}{p{0.3\linewidth}}{
        \cellcolor[rgb]{0.9,0.9,0.9}{
            \makecell[{{p{\linewidth}}}]{
                \texttt{\tiny{[GM$|$GM]}}
                \texttt{Jā, lūdzu! Ja tu esi pārliecināts par šo sadalījumu, tu vari nosūtīt noslēpto priekšlikumu. Es arī esmu gatavs to izdarīt, lai mēs sasniegtu šo rezultātu.} \\
            }
        }
    }
    & & \\ \\

    \theutterance \stepcounter{utterance}  
    & & \multicolumn{4}{p{0.6\linewidth}}{
        \cellcolor[rgb]{0.9,0.9,0.9}{
            \makecell[{{p{\linewidth}}}]{
                \texttt{\tiny{[GM$\rangle$P2]}}
                \texttt{Jā, lūdzu! Ja tu esi pārliecināts par šo sadalījumu, tu vari nosūtīt noslēpto priekšlikumu. Es arī esmu gatavs to izdarīt, lai mēs sasniegtu šo rezultātu.} \\
            }
        }
    }
    & \\ \\

    \theutterance \stepcounter{utterance}  
    & & & \multicolumn{4}{p{0.6\linewidth}}{
        \cellcolor[rgb]{0.9,0.9,0.9}{
            \makecell[{{p{\linewidth}}}]{
                \texttt{\tiny{[GM$\langle$P2]}}
                \texttt{{[}Priekšlikums: 3 grāmatas, 1 soma{]}} \\
            }
        }
    }
     \\ \\

    \theutterance \stepcounter{utterance}  
    & & & \multicolumn{2}{p{0.3\linewidth}}{
        \cellcolor[rgb]{0.9,0.9,0.9}{
            \makecell[{{p{\linewidth}}}]{
                \texttt{\tiny{[GM$|$GM]}}
                \texttt{proposal} \\
            }
        }
    }
    & & \\ \\

    \theutterance \stepcounter{utterance}  
    & & & \multicolumn{2}{p{0.3\linewidth}}{
        \cellcolor[rgb]{0.9,0.9,0.9}{
            \makecell[{{p{\linewidth}}}]{
                \texttt{\tiny{[GM$|$GM]}}
                [3, 1, 0]
            }
        }
    }
    & & \\ \\

    \theutterance \stepcounter{utterance}  
    & & \multicolumn{4}{p{0.6\linewidth}}{
        \cellcolor[rgb]{0.9,0.9,0.9}{
            \makecell[{{p{\linewidth}}}]{
                \texttt{\tiny{[P1$\langle$GM]}}
                \texttt{Otrs spēlētājs ir iesniedzis slepenu priekšlikumu. Tagad ir jūsu kārta iesniegt savu priekšlikumu.} \\
            }
        }
    }
    & \\ \\

    \theutterance \stepcounter{utterance}  
    & \multicolumn{4}{p{0.6\linewidth}}{
        \cellcolor[rgb]{0.9,0.9,0.9}{
            \makecell[{{p{\linewidth}}}]{
                \texttt{\tiny{[P1$\rangle$GM]}}
                \texttt{{[}Priekšlikums: 1 bultskrūve{]}} \\
            }
        }
    }
    & & \\ \\

    \theutterance \stepcounter{utterance}  
    & & & \multicolumn{2}{p{0.3\linewidth}}{
        \cellcolor[rgb]{0.9,0.9,0.9}{
            \makecell[{{p{\linewidth}}}]{
                \texttt{\tiny{[GM$|$GM]}}
                \texttt{proposal} \\
            }
        }
    }
    & & \\ \\

    \theutterance \stepcounter{utterance}  
    & & & \multicolumn{2}{p{0.3\linewidth}}{
        \cellcolor[rgb]{0.9,0.9,0.9}{
            \makecell[{{p{\linewidth}}}]{
                \texttt{\tiny{[GM$|$GM]}}
                [0, 0, 1]
            }
        }
    }
    & & \\ \\

    \theutterance \stepcounter{utterance}  
    & & & \multicolumn{2}{p{0.3\linewidth}}{
        \cellcolor[rgb]{0.9,0.9,0.9}{
            \makecell[{{p{\linewidth}}}]{
                \texttt{\tiny{[GM$|$GM]}}
                [[0, 0, 1], [3, 1, 0]]
            }
        }
    }
    & & \\ \\

\end{supertabular}
}

\end{document}
